\section{Instrumental variables}
Write
\[
Q=E[z_i z_i'],\qquad D=E[z_i x_i'],\qquad
S=E[\varepsilon_i^2z_i z_i'].
\]
The numbers of regressors and instruments are fixed. In addition to the
assumed i.i.d.\ sampling, impose sufficient moments for the laws of large
numbers and central limit theorem below; finite fourth moments of
$x_i,z_i,\varepsilon_i$ suffice. For (c)--(d), assume $S$ is positive definite.

\subsection*{(a) Valid instruments and consistency}
The instruments require \emph{orthogonality}, $E[z_i\varepsilon_i]=0$,
and \emph{relevance}, $\operatorname{rank}(D)=k$.
Also require $Q$ to be positive definite, removing redundant instruments.
The stronger condition $E[\varepsilon_i\mid z_i]=0$ is sufficient for
orthogonality but is not necessary.

Let
\[
Q_n=\frac{Z'Z}{n},\qquad D_n=\frac{Z'X}{n},\qquad
g_n=\frac{Z'\varepsilon}{n}.
\]
Since $\widehat W=Z(Z'Z)^{-1}Z'X=P_ZX$, where
$P_Z=Z(Z'Z)^{-1}Z'$ is symmetric and idempotent,
\[
\widehat\beta_{2SLS}
=(X'P_ZX)^{-1}X'P_ZY.
\]
Substituting $Y=X\beta+\varepsilon$ gives
\[
\widehat\beta_{2SLS}-\beta
=(D_n'Q_n^{-1}D_n)^{-1}D_n'Q_n^{-1}g_n.
\]
By the law of large numbers,
$Q_n\xrightarrow{p}Q$, $D_n\xrightarrow{p}D$ and
$g_n\xrightarrow{p}E[z_i\varepsilon_i]=0$.
Since $Q$ is positive definite and $D$ has full column rank,
$D'Q^{-1}D$ is positive definite. Thus
\[
\boxed{\widehat\beta_{2SLS}\xrightarrow{p}\beta.}
\]

\subsection*{(b) Asymptotic variance under heteroskedasticity}
Let $H=D'Q^{-1}D$. Since orthogonality makes $E[z_i\varepsilon_i]=0$,
the multivariate central limit theorem gives
\[
\frac1{\sqrt n}\sum_{i=1}^n z_i\varepsilon_i
\xrightarrow{d}N(0,S),\qquad
S=E[\sigma_i^2z_i z_i'],\quad
\sigma_i^2=E[\varepsilon_i^2\mid z_i].
\]
Consequently,
\[
\sqrt n(\widehat\beta_{2SLS}-\beta)
=
H^{-1}D'Q^{-1}
\frac1{\sqrt n}\sum_{i=1}^n z_i\varepsilon_i+o_p(1),
\]
so
\[
\boxed{
\sqrt n(\widehat\beta_{2SLS}-\beta)\xrightarrow{d}N(0,V_{2SLS}),
\quad
V_{2SLS}=H^{-1}D'Q^{-1}SQ^{-1}DH^{-1}.
}
\]
This is the covariance of the limiting distribution; the large-sample
covariance of $\widehat\beta_{2SLS}$ itself is $V_{2SLS}/n$.

\subsection*{(c) Residual-based minimum chi-square estimation}
The matrix defined in the question satisfies
\[
\widehat S=\frac{\widehat R}{n}
=\frac1n\sum_{i=1}^n z_i z_i'\widehat\varepsilon_i^2,
\qquad
\widehat\varepsilon_i
=\varepsilon_i-x_i'(\widehat\beta_{2SLS}-\beta).
\]
Expanding the square gives
\[
\widehat S-\frac1n\sum_i z_i z_i'\varepsilon_i^2=o_p(1).
\]
Indeed, the difference contains terms linear and quadratic in the
consistent estimation error $\widehat\beta_{2SLS}-\beta$,
multiplied by sample averages that are $O_p(1)$ under the moment
conditions. Hence $\widehat S\xrightarrow{p}S$.

Holding this first-step residual weight fixed, minimization of the
criterion in the question yields
\[
\widehat\beta_{MCS}
=(X'Z\widehat R^{-1}Z'X)^{-1}X'Z\widehat R^{-1}Z'Y.
\]
Since $\widehat R=n\widehat S$,
\[
\widehat\beta_{MCS}-\beta
=(D_n'\widehat S^{-1}D_n)^{-1}
 D_n'\widehat S^{-1}g_n.
\]
The coefficient matrix converges to a finite matrix and $g_n\to_p0$,
so $\boxed{\widehat\beta_{MCS}\xrightarrow{p}\beta}$.
Furthermore,
\[
\sqrt n(\widehat\beta_{MCS}-\beta)
=(D'S^{-1}D)^{-1}D'S^{-1}
\frac1{\sqrt n}\sum_{i=1}^n z_i\varepsilon_i+o_p(1).
\]
Thus
\[
\boxed{
\sqrt n(\widehat\beta_{MCS}-\beta)\xrightarrow{d}N(0,V_{MCS}),
\qquad V_{MCS}=(D'S^{-1}D)^{-1}.
}
\]

\subsection*{(d) Efficiency comparison}
Define the two asymptotic influence matrices
\[
B_Q=(D'Q^{-1}D)^{-1}D'Q^{-1},\qquad
B_*=(D'S^{-1}D)^{-1}D'S^{-1}.
\]
Both satisfy $B_QD=B_*D=I_k$. Let $\Delta=B_Q-B_*$, so $\Delta D=0$.
Also $SB_*'=D(D'S^{-1}D)^{-1}$, which implies
$\Delta SB_*'=0$. Expanding $B_Q=B_*+\Delta$ therefore gives
\[
\begin{aligned}
V_{2SLS}-V_{MCS}
&=B_QSB_Q'-B_*SB_*'\\
&=\Delta S\Delta'\succeq0.
\end{aligned}
\]
For every $a\in\mathbb R^k$,
$a'V_{MCS}a\leq a'V_{2SLS}a$: MCS is asymptotically at least as
efficient as 2SLS. The dominance is weak; heteroskedasticity alone
does not guarantee strict improvement. Equality holds when
$B_Q=B_*$, including the homoskedastic case $S=\sigma^2Q$.
